\documentclass[11pt,a4paper]{article}

% Required packages
\usepackage{graphicx}
\usepackage{float}
\usepackage{caption}
\usepackage{adjustbox}

\captionsetup{
    font=small,
    labelfont=bf
}

\usepackage[margin=1.5cm]{geometry}
\usepackage{pdflscape}
\usepackage{amsmath}
\usepackage{amsfonts}

\title{Table of Contents\\Chapters 1--4\\[0.3cm]\small{Randomized Reward Shaping Deep Reinforcement Learning\\for Self-Adaptive Systems}}
\author{Alireza Kavianifar}
\date{\today}

\begin{document}

\maketitle

\tableofcontents
\newpage

% ============================================================================
% CHAPTER 1: INTRODUCTION
% ============================================================================

% Set chapter counter to 1 so sections are numbered 1.1, 1.2, etc.
\setcounter{section}{0}
\renewcommand{\thesection}{1.\arabic{section}}
\renewcommand{\thesubsection}{1.\arabic{section}.\arabic{subsection}}
\renewcommand{\thesubsubsection}{1.\arabic{section}.\arabic{subsection}.\arabic{subsubsection}}

\section{Background and Context}
\subsection{Self-Adaptive Systems: motivation and examples}
\subsection{The MAPE-K feedback loop and the Managed/Managing System abstraction}
\subsection{Limitations of rule-based and static adaptation}

\section{Runtime Uncertainty in Self-Adaptive Systems}
\subsection{Sources of uncertainty (environmental, system-level)}
\subsection{Impact of uncertainty on adaptation decisions (examples from DeltaIoT)}

\section{Learning-based Adaptation: Why Reinforcement Learning?}
\subsection{RL benefits for runtime decision-making}
\subsection{Limits of classic RL in large adaptation spaces}

\section{What ``Continuous'' Means in This Thesis}
\subsection{Continuous = continuous runtime learning / streaming adaptation cycles (not continuous-action control)}

\section{Problem Statement (precise formulation extracted from the article)}

\section{Proposed Solution Overview: RS-DRL (high-level)}
\subsection{Key idea: randomized reward reshaping during experience replay}
\subsection{Architectural placement within MAPE-K (on-demand retraining \& async learning)}

\section{Research Objectives and Research Questions}

\section{Contributions of the Thesis (expanded from the article's contribution list)}

\section{Research Methodology Overview}
\subsection{Theoretical modeling and RL formulation}
\subsection{Algorithm design and reward shaping experiments}
\subsection{Case study (DeltaIoT): datasets, training regime, evaluation metrics}

\section{Thesis Organization (brief outline of subsequent chapters)}

\newpage

% ============================================================================
% CHAPTER 2: FOUNDATIONS AND RELATED WORK
% ============================================================================

\setcounter{section}{0}
\renewcommand{\thesection}{2.\arabic{section}}
\renewcommand{\thesubsection}{2.\arabic{section}.\arabic{subsection}}
\renewcommand{\thesubsubsection}{2.\arabic{section}.\arabic{subsection}.\arabic{subsubsection}}

\section{Self-Adaptive Systems: Concepts and Architectural Patterns}
\subsection{MAPE-K detailed (Monitor, Analyze, Plan, Execute, Knowledge)}
\subsection{Managed vs Managing System abstraction and temporal classification (monitoring vs asynchronous learning)}

\section{Uncertainty in Self-Adaptive Systems}
\subsection{Taxonomy of uncertainties (environmental, systemic, goal drift)}
\subsection{Effects of uncertainty on verification and adaptation (motivating evidence from DeltaIoT)}

\section{Reinforcement Learning Foundations for SAS}
\subsection{MDPs, policy/value formulation and why RL suits adaptation problems}
\subsection{Known challenges: local optima, generalization, retraining cost}

\section{Deep Reinforcement Learning in Large Discrete Adaptation Spaces}
\subsection{DQN architecture and rationale for discrete adaptation options (why DQN was chosen in the article)}
\subsection{Experience replay, target networks, stability techniques used}

\section{Reward Shaping and Experience Relabeling (state of the art)}
\subsection{HER and goal-conditioned methods --- contrast to RS-DRL (no pre-defined goals)}
\subsection{Other generalization techniques}

\section{Related Work: classification and gaps}
\subsection{ML-based methods, search-based planners, RL-based adaptation --- strengths and weaknesses (article's survey)}
\subsection{Identified gap: scalable, uncertainty-aware, continuous runtime DRL for large adaptation spaces}

\section{Chapter Summary (lead into Chapter 3's formal modeling)}

\newpage

% ============================================================================
% CHAPTER 3: SYSTEM MODEL AND PROBLEM FORMULATION
% ============================================================================

\setcounter{section}{0}
\renewcommand{\thesection}{3.\arabic{section}}
\renewcommand{\thesubsection}{3.\arabic{section}.\arabic{subsection}}
\renewcommand{\thesubsubsection}{3.\arabic{section}.\arabic{subsection}.\arabic{subsubsection}}

\section{Motivating Case Study: DeltaIoT System (detailed description)}
\subsection{Network, motes, links, configuration parameters and examples}

\section{Modeling Uncertainty in DeltaIoT and SAS}
\subsection{Interference/noise ranges, message-load variability, data shifts (Fig.2 evidence)}
\subsection{How uncertainty affects QoS metrics (packet loss, latency, energy)}

\section{Adaptation Space and Configuration Representation}
\subsection{Definition: adaptation option, parameter encoding, scale (216 / 1,296 / 4,096)}
\subsection{Implications for search/learning complexity}

\section{RL Formalization for the Adaptive Problem (exact equations from paper)}
\subsection{State space definition (e.g., $S = (E, P, L)$)}
\subsection{Action space: discrete adaptation options and encoding}
\subsection{Reward functions: multi-objective normalization, threshold/setpoint formulations (include equations 11--12 and Table 2)}
\subsection{Transition dynamics (how verifier / ActivFORMS produces next states)}

\section{Continuous / Streaming Training Assumptions}
\subsection{Treating adaptation cycles as a continuous stream (training horizons and validation splits per instance)}
\subsection{Temporal separation: runtime adaptation vs asynchronous learning (design assumption)}

\section{Verification \& Runtime Models for Option Analysis}
\subsection{UPPAAL-SMC models and ActivFORMS execution engine (role in reward calculation)}

\section{Formal Problem Statement (thesis version) --- precise objective and constraints (ready to feed Chapter 4)}

\newpage

% ============================================================================
% CHAPTER 4: ARCHITECTURE AND RS-DRL METHOD
% ============================================================================

\setcounter{section}{0}
\renewcommand{\thesection}{4.\arabic{section}}
\renewcommand{\thesubsection}{4.\arabic{section}.\arabic{subsection}}
\renewcommand{\thesubsubsection}{4.\arabic{section}.\arabic{subsection}.\arabic{subsubsection}}

\section{Chapter overview \& architectural perspective (scope, process-oriented organization)}

\section{System-level Self-Adaptation Loop (nine-stage continuous adaptation cycle)}
\subsection{Stage 1: Runtime data collection (sensors $\rightarrow$ Monitor)}
\subsection{Stage 2: State forwarding to RS-DRL module (experience collection)}
\subsection{Stage 3: Asynchronous model training and storage (replay $\rightarrow$ train $\rightarrow$ store)}
\subsection{Stage 4: State forwarding to Analyzer}
\subsection{Stage 5: Adaptation triggering (Analyzer detects violations)}
\subsection{Stage 6: Trained model retrieval}
\subsection{Stage 7: Adaptation option prediction (Adaptation Option Predictor)}
\subsection{Stage 8: Adaptation option retrieval \& verification}
\subsection{Stage 9: Adaptation application (Executor)}

\section{Mapping RS-DRL to MAPE-K (conceptual mapping and roles)}

\section{Runtime Monitoring and State Construction}
\subsection{Data collection, preprocessing, feature extraction, state vector construction}

\section{Analysis and Adaptation Triggering}
\subsection{Change detection \& thresholding assumptions (when retraining is triggered)}

\section{Planning \& Adaptation Option Prediction}
\subsection{Adaptation Option Selector ($\epsilon$-greedy + DQN) \& architecture of predictor}

\section{Verification of Adaptation Options}
\subsection{Runtime model checking with UPPAAL-SMC \& how verifier informs reward}

\section{Execution and Adaptation Application}
\subsection{Executor behavior and coordination with Knowledge repository}

\section{Asynchronous RS-DRL Learning and Policy Update (core algorithmic chapter)}
\subsection{Overall RS-DRL algorithm (Algorithm 1) and parameters ($\eta$, $\gamma$, $\epsilon$, $\rho$)}
\subsection{Replay Memory \& randomized reward reshaping (Algorithm 2) --- rationale \& implementation}
\subsection{Model training, target networks, loss function, stability practices}

\section{Knowledge Management and Process Coordination}
\subsection{Knowledge repository design (models, adaptation options, state history)}
\subsection{Read/write coordination contracts and heterogeneous process synchronization}

\section{Implementation details \& hyperparameter choices}
\subsection{Gymnasium environment, Stable-Baselines3 DQN extension, dataset traces (Gheibi \& Weyns)}

\section{Chapter summary and transition to evaluation (Chapter 5)}

\end{document}
